Une bouteille contient un volume $V=\qty{230}{\cubic\cm}$ du dioxygène
\ce{O2} gazeux sous la pression $P=\qty{1033}{\hPa}$
et à la température $\theta=\qty{25}{\degreeCelsius}$.
\begin{enumerate}
    %\item Déterminer la densité du dioxygène par rapport à l'air.
    \item Calculer la quantité de matière du gaz dioxygène qui se
          trouve dans cette bouteille (en le considérant comme un gaz
          parfait).
    \item Déterminer la valeur du volume molaire dans les conditions
          précédentes.
    \item Quelle est la pression qu'on doit exercer sur l'échantillon
          du gaz précédent à la température
          $\theta^{\prime}=\qty{20}{\degreeCelsius}$ pour que son
          volume devient $V^{\prime}=\qty{0,8}{\liter}$.
\end{enumerate}
\begin{donnees}
       $M(\ce{O2})=\qty{32}{\gram\per\mol}$, \
       $R=\qty{8,314}{\J\per\mol\per\K}$, \
       $\qty{1}{\liter}=\qty{1e-3}{\cubic\meter}$, \
       $\qty{1}{\hecto\Pa}=\qty{100}{\Pa}$.
\end{donnees}
